\documentclass[aspectratio=169, 10pt]{beamer}
\usepackage{beamer_theme_cienciadados2026}

\title[Aula 02 - Ecossistema Jupyter]{Ciência dos Dados I: Análise Exploratória}
\subtitle{Aula 02: O Ecossistema Jupyter e Computação Científica Interativa}
\author{Prof. Dr. Ney Lemke}
\institute{UNESP -- Instituto de Biociências de Botucatu}
\date{Versão 2026}

\begin{document}

\frame{\titlepage}

\begin{frame}{Sumário da Aula}
  \tableofcontents
\end{frame}

\section{Origem e Filosofia do Jupyter}

\begin{frame}{Do IPython ao Projeto Jupyter}
  \begin{itemize}
    \item \textbf{2001}: Fernando Pérez cria o \textbf{IPython} para trazer interatividade superior ao interpretador Python padrão.
    \item \textbf{2011}: Lançamento do formato de notebook baseado em navegador web (\texttt{.ipynb}).
    \item \textbf{2014}: Nascimento do \textbf{Projeto Jupyter} (acrônimo para \textbf{Ju}lia, \textbf{Pyt}hon e \textbf{R}), desacoplando a interface da linguagem de programação.
    \item \textbf{Hoje (2026)}: Padrão universal na academia e indústria (JupyterLab 4+, Google Colab, VS Code, Kaggle).
  \end{itemize}
\end{frame}

\section{Arquitetura do Jupyter}

\begin{frame}{Arquitetura Cliente-Servidor e Protocolo ZeroMQ}
  \begin{columns}
    \begin{column}{0.5\textwidth}
      \begin{block}{1. Interface do Usuário (Cliente)}
        Navegador web enviando comandos e renderizando saídas (HTML, SVG, Canvas, LaTeX, imagens).
      \end{block}
      \begin{block}{2. Servidor Jupyter}
        Gerencia autenticação, arquivos no disco e conexões de rede via WebSockets.
      \end{block}
    \end{column}
    \begin{column}{0.5\textwidth}
      \begin{block}{3. Kernel (IPykernel)}
        Processo isolado que executa o código Python na memória e retorna os resultados via mensagens estruturadas JSON (ZeroMQ).
      \end{block}
      \begin{block}{Vantagem Arquitetural}
        O kernel pode estar rodando em uma GPU remota na nuvem enquanto a interface roda localmente no seu notebook.
      \end{block}
    \end{column}
  \end{columns}
\end{frame}

\section{Reprodutibilidade Científica}

\begin{frame}{O Desafio da Reprodutibilidade em Notebooks}
  \begin{alertblock}{O Perigo do Estado Oculto}
    Notebooks permitem execução não-linear de células. Se você executar a célula 5, depois a célula 2 e depois a 4, o estado da memória pode se tornar inconsistente e irreprodutível para outros pesquisadores!
  \end{alertblock}
  \begin{block}{Regra de Ouro do Cientista de Dados}
    Antes de submeter um relatório ou compartilhar um notebook, execute sempre: \\
    \textbf{Kernel $\rightarrow$ Restart Kernel and Run All Cells} \\
    Garanta que o caderno execute do início ao fim sem falhas.
  \end{block}
\end{frame}

\section{Recursos Avançados: Markdown e Magics}

\begin{frame}[fragile]{Comandos Mágicos do IPython}
  Comandos especiais que estendem a capacidade do Python:
  \begin{itemize}
    \item \texttt{\%timeit}: Mede o tempo médio de execução de instruções com alta precisão estatística.
    \item \texttt{\%whos}: Lista todas as variáveis ativas em memória com seus tipos e tamanhos.
    \item \texttt{!\textit{comando}}: Executa comandos no shell do sistema operacional diretamente da célula.
  \end{itemize}
  \begin{lstlisting}
# Exemplo no notebook:
%timeit sum(range(1_000_000))
!ls -lh dados/
  \end{lstlisting}
\end{frame}

\begin{frame}{Resumo da Aula}
  \begin{itemize}
    \item Jupyter combina narrativa, código e visualização em documentos executáveis vivos.
    \item O kernel é o motor de execução e deve ser mantido limpo e reprodutível.
    \item \textbf{Prática}: Dominar comandos mágicos, Markdown técnico e equações LaTeX.
  \end{itemize}
\end{frame}

\end{document}
